% SolarMax — Flashing the ESP32 (VS Code + PlatformIO)
% Compile:  pdflatex flashing_guide.tex   (run twice for the table of contents)
% Real screenshots use \realshot{...}; missing (hardware) shots stay as \shot{...} placeholders.

\documentclass[11pt]{article}
\usepackage[margin=2.2cm]{geometry}
\usepackage{graphicx}
\usepackage{float}
\usepackage{xcolor}
\usepackage[most]{tcolorbox}
\usepackage{enumitem}
\usepackage{titlesec}
\usepackage[colorlinks=true,linkcolor=blue!60!black,urlcolor=blue!60!black]{hyperref}

\graphicspath{{./}}

% ── Repo location ─────────────────────────────────────────────────────────────
\newcommand{\repourl}{\href{https://github.com/zhsquared/SolarMax}{github.com/zhsquared/SolarMax}}

% ── Colors / section styling ─────────────────────────────────────────────────
\definecolor{brandblue}{HTML}{1565C0}
\definecolor{brandgreen}{HTML}{2E7D32}
\titleformat{\section}{\Large\bfseries\color{brandblue}}{\thesection.}{0.5em}{}

% ── Callout boxes ─────────────────────────────────────────────────────────────
\newtcolorbox{tip}{colback=green!6,colframe=brandgreen,boxrule=0.8pt,
  left=8pt,right=8pt,top=6pt,bottom=6pt,title=\textbf{Tip},fonttitle=\color{white}}
\newtcolorbox{warn}{colback=red!5,colframe=red!70!black,boxrule=0.8pt,
  left=8pt,right=8pt,top=6pt,bottom=6pt,title=\textbf{Important},fonttitle=\color{white}}
\newtcolorbox{note}{colback=blue!5,colframe=brandblue,boxrule=0.8pt,
  left=8pt,right=8pt,top=6pt,bottom=6pt,title=\textbf{Note},fonttitle=\color{white}}

% ── A real screenshot: #1 = image file, #2 = caption ─────────────────────────
\newcommand{\realshot}[2]{%
\begin{center}
\fbox{\includegraphics[width=0.86\linewidth]{#1}}\\[3pt]
{\small\itshape #2}
\end{center}}

% ── A still-needed screenshot: #1 = intended filename, #2 = what to capture ──
\newcommand{\shot}[2]{%
\begin{center}
\setlength{\fboxsep}{10pt}%
\fcolorbox{gray}{gray!8}{\parbox[c][2.8cm][c]{0.86\linewidth}{\centering
\textcolor{gray!55!black}{\small\ttfamily [\,IMAGE PLACEHOLDER:\ #1\,]}\\[6pt]
\textbf{Screenshot to take:}\ #2}}
\end{center}}

\setlist[enumerate]{leftmargin=1.6em,itemsep=3pt,topsep=4pt}

\begin{document}

% ── Title ─────────────────────────────────────────────────────────────────────
\begin{center}
{\Huge\bfseries\color{brandblue} SolarMax}\\[4pt]
{\LARGE Flashing the ESP32 --- Step-by-Step Guide}\\[8pt]
{\large VS Code + PlatformIO. A beginner-friendly walkthrough --- no prior setup needed.}\\[6pt]
\rule{0.9\linewidth}{0.8pt}
\end{center}

\vspace{4pt}
\tableofcontents
\vspace{8pt}

% ── Before you start ──────────────────────────────────────────────────────────
\section{Before you start}

You need these on hand:
\begin{itemize}[leftmargin=1.6em]
  \item A computer (Windows or macOS) with internet.
  \item The ESP32 DevKit V1 board.
  \item A \textbf{USB data cable} (Micro-USB) that fits the board.
\end{itemize}

\begin{warn}
Many USB cables are \textbf{charge-only} and carry no data --- the board will power on
but the computer won't see it. If your board is not detected later, the cable is the
\#1 suspect. Use a known data cable (e.g.\ one that syncs a phone).
\end{warn}

\shot{board-overview.jpg}{A clear photo of the ESP32 board, with the USB port, the
\texttt{EN} (reset) button, and the \texttt{BOOT} button visible.}

% =============================================================================
\section{Install Visual Studio Code}
\begin{enumerate}
  \item Go to \href{https://code.visualstudio.com}{code.visualstudio.com}.
  \item Click the big download button for your system (Windows or Mac).
  \item Open the downloaded file and follow the installer (accept the defaults).
  \item Launch VS Code.
\end{enumerate}
\realshot{vscode-01-download.png}{The Visual Studio Code download page. Click the large
``Download for macOS'' button (or the \textbf{Download} button at the top-right for your
operating system).}

% =============================================================================
\section{Install the PlatformIO extension}
\begin{enumerate}
  \item In VS Code, click the \textbf{Extensions} icon in the left bar (four squares).
  \item In the search box, type \texttt{PlatformIO IDE}.
  \item Click \textbf{Install} on the ``PlatformIO IDE'' result.
  \item Wait for it to finish --- this can take a few minutes. When prompted, allow it
        to reload/restart VS Code.
\end{enumerate}
\realshot{vscode-02-platformio-extension.png}{The PlatformIO IDE extension page. Click
\textbf{Install}. (Here it is already installed, so it shows Disable/Uninstall instead.)}

\begin{tip}
When PlatformIO finishes installing, a small \textbf{alien-head icon} appears in the
left bar. Clicking it opens \textbf{PlatformIO Home} --- that's how you know the install
worked.
\end{tip}
\realshot{vscode-03-pio-home.png}{PlatformIO Home. The alien-head icon in the far-left
bar (highlighted near the bottom) opens it. From here you can ``Pick a folder'' to open a
project.}

% =============================================================================
\section{Get the SolarMax project onto the computer}
Pick \emph{one} option:
\begin{itemize}[leftmargin=1.6em]
  \item \textbf{Easiest --- download a copy:} open \repourl{} in a browser, click the
        green \textbf{Code} button, choose \textbf{Download ZIP}, then unzip it. You'll
        get a folder named \texttt{SolarMax}.
  \item \textbf{Or copy it:} get the \texttt{SolarMax} folder from a teammate (USB drive
        or shared drive).
\end{itemize}
\realshot{vscode-04-download-zip.png}{On the GitHub repo page, click the green
\textbf{Code} button and choose \textbf{Download ZIP} (or copy the HTTPS clone URL shown).}

% =============================================================================
\section{Open the project in VS Code}
\begin{enumerate}
  \item In VS Code: \textbf{File \(\rightarrow\) Open Folder\ldots}
  \item Select the \texttt{SolarMax} folder (the one that contains
        \texttt{platformio.ini}) and open it.
  \item Wait. The first time, PlatformIO downloads the ESP32 tools automatically --- a
        progress bar shows at the bottom. Let it finish before continuing.
\end{enumerate}
\realshot{vscode-05-open-folder.png}{The SolarMax project open in VS Code. The Explorer on
the left lists \texttt{platformio.ini} plus the \texttt{src}, \texttt{lib},
\texttt{include}, \texttt{docs}, and \texttt{sim} folders.}

% =============================================================================
\section{Install the USB driver (only if the board isn't detected)}
Most computers already have this. If, during upload, no port is found, install the driver
for your board's USB chip:
\begin{itemize}[leftmargin=1.6em]
  \item \textbf{CP210x} (most DevKit V1 boards):
        \href{https://www.silabs.com/developer-tools/usb-to-uart-bridge-vcp-drivers}{Silicon Labs CP210x VCP driver}.
  \item \textbf{CH340} (some clones): search ``CH340 driver'' for your OS.
\end{itemize}
\realshot{driver-cp210x.png}{The Silicon Labs CP210x driver downloads. Choose the one for
your operating system (e.g.\ ``CP210x VCP Mac OSX Driver'', or a Windows driver).}

% =============================================================================
\section{Plug in the ESP32}
Connect the board to the computer with the USB \emph{data} cable. A power LED on the
board should light up.
\shot{board-plugged-in.jpg}{The ESP32 plugged into the computer, with its power LED lit.}

% =============================================================================
\section{Upload (flash) the firmware}
\begin{enumerate}
  \item Look at the \textbf{toolbar along the very bottom} of the VS Code window --- a
        thin dark bar with small icons (this is the PlatformIO toolbar).
  \item Click the \textbf{right-arrow icon} (\(\rightarrow\)) --- ``PlatformIO: Upload''.
        (The checkmark next to it is Build-only; the arrow both builds \emph{and} uploads.)
  \item The terminal at the bottom compiles the code and flashes the board. Wait for it.
  \item Success looks like \texttt{[SUCCESS]} and ``Took \ldots seconds'' in green.
\end{enumerate}
\realshot{vscode-06-toolbar.png}{The PlatformIO toolbar at the bottom of the window. Left
to right: home, checkmark (\textbf{Build}), right-arrow (\textbf{Upload}), and further
right the plug icon (\textbf{Serial Monitor}). ``Default (SolarMax)'' shows the active
project.}
\shot{vscode-07-upload-success.png}{The VS Code terminal showing a successful upload
ending in \texttt{[SUCCESS]}.}

\begin{warn}
If upload fails with a ``Connecting\ldots'' timeout: hold the board's \textbf{BOOT}
button while it says ``Connecting'', then release it once it starts writing. Also
double-check the cable (Section~1) and the driver (Section~6).
\end{warn}

% =============================================================================
\section{Watch it run (Serial Monitor)}
\begin{enumerate}
  \item In the same bottom toolbar, click the \textbf{plug icon} (``Serial Monitor'').
  \item You should see the SolarMax start-up banner and the self-check report.
\end{enumerate}
\shot{vscode-08-serial-monitor.png}{The VS Code Serial Monitor showing SolarMax boot
output (the ``SolarMax'' banner and the self-check lines).}

\begin{tip}
That's it --- the ESP32 is flashed. To flash again after code changes, just click the
Upload arrow again.
\end{tip}

% =============================================================================
\section{Troubleshooting}
\begin{itemize}[leftmargin=1.6em]
  \item \textbf{No port / board not detected:} try a different (data) USB cable; install
        the USB driver (Section~6); try a different USB port.
  \item \textbf{Stuck on ``Connecting\ldots'':} hold \textbf{BOOT} during the connect
        step, release when writing starts.
  \item \textbf{Nothing on Serial Monitor:} make sure the baud rate is \textbf{115200}
        and the correct port is selected.
  \item \textbf{Upload works but behaves oddly:} press the board's \textbf{EN} (reset)
        button once to restart it.
\end{itemize}

\end{document}
